% =====================================================================
%  manual.tex -- SparkRoad-V 数字示波器 产品手册
%  编译: xelatex manual.tex  (两次以生成目录)
% =====================================================================
\documentclass[UTF8,a4paper,11pt]{ctexart}
% =====================================================================
%  style.tex -- 公共宏包、版式与代码排版设置
% =====================================================================
\usepackage[a4paper,left=2.0cm,right=2.0cm,top=2.4cm,bottom=2.2cm,headsep=0.6cm]{geometry}
\usepackage{xcolor}
\usepackage{listings}
\usepackage{booktabs}
\usepackage{longtable}
\usepackage{array}
\usepackage{amsmath,amssymb}
\usepackage{caption}
\usepackage{enumitem}
\usepackage{fancyhdr}
\usepackage{tikz}
\usetikzlibrary{positioning,arrows.meta,shapes.geometric,fit,backgrounds}
\usepackage[colorlinks=true,linkcolor=blue!55!black,urlcolor=blue!55!black,citecolor=blue!55!black,
            bookmarksnumbered=true]{hyperref}

% ---------- 颜色 ----------
\definecolor{codebg}{gray}{0.985}
\definecolor{codeframe}{gray}{0.62}
\definecolor{kw}{RGB}{0,80,160}
\definecolor{cmt}{RGB}{0,110,90}
\definecolor{str}{RGB}{150,60,0}
\definecolor{warnbg}{RGB}{255,244,224}
\definecolor{warnframe}{RGB}{200,120,0}
\definecolor{tipbg}{RGB}{232,244,255}
\definecolor{tipframe}{RGB}{30,90,170}

% ---------- 代码块样式 ----------
\lstdefinestyle{vcode}{
  language=Verilog,
  basicstyle=\ttfamily\scriptsize,
  keywordstyle=\color{kw}\bfseries,
  commentstyle=\color{cmt},
  stringstyle=\color{str},
  numbers=left, numberstyle=\tiny\color{gray}, numbersep=6pt,
  backgroundcolor=\color{codebg}, frame=single, rulecolor=\color{codeframe},
  framexleftmargin=1.2em, xleftmargin=1.6em,
  columns=fullflexible, keepspaces=true, showstringspaces=false,
  tabsize=4, breaklines=true, breakatwhitespace=false,
  extendedchars=true, inputencoding=utf8,
  aboveskip=1.0em, belowskip=0.6em,
  upquote=true
}
\lstdefinestyle{vsnip}{
  language=Verilog,
  basicstyle=\ttfamily\small,
  keywordstyle=\color{kw}\bfseries,
  commentstyle=\color{cmt},
  stringstyle=\color{str},
  numbers=none,
  backgroundcolor=\color{codebg}, frame=single, rulecolor=\color{codeframe},
  framexleftmargin=0.6em, xleftmargin=1.0em,
  columns=fullflexible, keepspaces=true, showstringspaces=false,
  tabsize=4, breaklines=true, breakatwhitespace=false,
  extendedchars=true, inputencoding=utf8,
  aboveskip=0.8em, belowskip=0.5em
}
\lstdefinestyle{cfgcode}{
  basicstyle=\ttfamily\scriptsize,
  numbers=left, numberstyle=\tiny\color{gray},
  backgroundcolor=\color{codebg}, frame=single, rulecolor=\color{codeframe},
  framexleftmargin=1.2em, xleftmargin=1.6em,
  columns=fullflexible, keepspaces=true, showstringspaces=false,
  breaklines=true, extendedchars=true, inputencoding=utf8
}

% ---------- 标识符排版 ----------
\DeclareRobustCommand{\sig}[1]{\texttt{\detokenize{#1}}}
\DeclareRobustCommand{\file}[1]{\texttt{\detokenize{#1}}}
\DeclareRobustCommand{\reg}[1]{\sig{#1}}

% ---------- 提示框(环境形式, 内部可用段落/公式/列表) ----------
\newsavebox{\noteboxsave}
\newenvironment{warnbox}
  {\par\medskip\noindent\begin{lrbox}{\noteboxsave}%
   \begin{minipage}{\dimexpr\linewidth-2\fboxsep-2\fboxrule}\small}
  {\end{minipage}\end{lrbox}%
   \fcolorbox{warnframe}{warnbg}{\usebox{\noteboxsave}}\par\smallskip}
\newenvironment{tipbox}
  {\par\medskip\noindent\begin{lrbox}{\noteboxsave}%
   \begin{minipage}{\dimexpr\linewidth-2\fboxsep-2\fboxrule}\small}
  {\end{minipage}\end{lrbox}%
   \fcolorbox{tipframe}{tipbg}{\usebox{\noteboxsave}}\par\smallskip}

% ---------- 表格辅助 ----------
\newcolumntype{L}[1]{>{\raggedright\arraybackslash}p{#1}}
\newcolumntype{C}[1]{>{\centering\arraybackslash}p{#1}}
\newcommand{\tabfont}{\small}
\captionsetup{font=small,labelfont=bf}
\setlength{\parskip}{0.25em}
\renewcommand{\arraystretch}{1.15}

% ---------- 页眉页脚 ----------
\pagestyle{fancy}
\fancyhf{}
\fancyhead[L]{\small 高云 GW5A-25A 的 \texttt{project\_2025\_11\_07} 工程代码逐层解析}
\fancyhead[R]{\small \thepage}
\renewcommand{\headrulewidth}{0.3pt}

% ---------- 目录深度 ----------
\setcounter{tocdepth}{3}
\setcounter{secnumdepth}{3}

% ---------- 常用宏 ----------
\newcommand{\cntmax}[1]{\sig{#1}}
\newcommand{\unit}[1]{\,\mathrm{#1}}


% ---- 手册专用页眉（覆盖主文档的设置）----
\pagestyle{fancy}
\fancyhf{}
\fancyhead[L]{\small SparkRoad-V 数字存储示波器}
\fancyhead[R]{\small 产品使用手册 \quad v1.0}
\fancyfoot[C]{\small \thepage}
\renewcommand{\headrulewidth}{0.4pt}
\renewcommand{\footrulewidth}{0pt}

\begin{document}

% ======================= 封面 =======================
\begin{titlepage}
\centering
\vspace*{2.6cm}
{\Huge\bfseries SparkRoad-V\\[3mm] 数字存储示波器}\\[6mm]
{\LARGE 产品使用手册}\\[16mm]

\begin{tabular}{rl}
\toprule
产品型号 & \sig{SPARKROAD-V-SCOPE}（EG4S20BG256 平台）\\
固件版本 & v1.0（P1：4 键操作 + 内置函数发生器）\\
显示输出 & VGA 640$\times$480 @ 60\,Hz\\
输入通道 & CH0 / CH1（AD7606，16\,bit，200\,kSa/s）\\
采样深度 & 1024 点 / 通道\\
内置信号源 & 正弦 / 方波 / 三角 / 锯齿，8 档频率，4 档幅度\\
操作方式 & 4 个功能按键 + 串口命令\\
文档版本 & 1.0\\
日期 & 2026 年 10 月\\
\bottomrule
\end{tabular}

\vfill
{\small 本手册面向\textbf{使用者}，只讲"怎么用"，不讲内部实现。\\
如需了解 RTL 设计细节，请参阅《FPGA 工程代码逐层解析与注释说明》。}
\end{titlepage}

\tableofcontents
\clearpage

% ======================= 1 产品简介 =======================
\section{产品简介}

本产品把一块 SparkRoad-V FPGA 开发板（安路 EG4S20BG256）做成了一台
\textbf{数字存储示波器}：采集——存储——触发——测量——显示全部在 FPGA 内部完成，
不依赖任何外部处理器。

\subsection{主要特性}

\begin{center}
\small
\begin{tabular}{ll}
\toprule
特性 & 说明\\
\midrule
双通道显示 & CH0（绿）/ CH1（黄），可单独关闭任一通道\\
VGA 直出 & 640$\times$480，行频 31.5\,kHz，场频 60.0\,Hz（片内 PLL 倍频）\\
实时采集 & AD7606 并行接口，200\,kSa/s，10\,bit 归一化后显示\\
触发系统 & 自动 / 正常 / 单次；上升沿 / 下降沿；电平可调；事件触发\\
预触发 & 可设置触发点之前的采样点数，能看到触发前波形\\
自动测量 & 峰峰值 Vpp、最大 Vmax、最小 Vmin、均值 Vavg、周期、频率、占空比\\
内置信号源 & 无需外接信号即可自测：4 种波形 / 8 档频率 / 4 档幅度\\
按键操作 & 4 键完成通道、垂直档位、时基、波形切换、菜单设置\\
串口控制 & 115200 8N1，命令式配置（详见第 \ref{sec:uart} 节）\\
\bottomrule
\end{tabular}
\end{center}

\subsection{前面板与接口}

\begin{center}
\small
\begin{tabular}{lll}
\toprule
接口 & 位置 & 说明\\
\midrule
VGA 输出 & 板载 VGA 座 & 接显示器，建议 640$\times$480 或更高分辨率\\
按键 L/R/U/D & 板载 4 个按键 & 见第 \ref{sec:key} 节\\
复位键 C & 板载 & 按下即整机复位（所有设置回到默认）\\
数码管 & 板载 4 位 & 显示 ADC 原始值\\
ADC 接口 & 2$\times$10 排针 & 接 AD7606 模块，未接时请使用内置信号源\\
USB-JTAG & 板载 & 烧写与调试\\
\bottomrule
\end{tabular}
\end{center}

\begin{tipbox}
\textbf{没有接 AD7606 也能测。}
内置函数发生器可替代外部信号，上电即有一条 50\,Hz 正弦波显示在屏幕上，
用于验证整条链路是否正常。
\end{tipbox}

% ======================= 2 显示界面 =======================
\section{显示界面说明}

\subsection{屏幕布局}

屏幕分为两块区域：

\begin{center}
\small
\begin{tabular}{lll}
\toprule
区域 & 像素范围 & 内容\\
\midrule
绘图区 & $y=0\ldots447$ & 网格、零轴、CH0 / CH1 波形\\
文字区 & $y=448\ldots463$ & 两行状态与测量读数\\
\bottomrule
\end{tabular}
\end{center}

\subsection{颜色约定}

\begin{center}
\small
\begin{tabular}{lll}
\toprule
元素 & 颜色 & 说明\\
\midrule
背景 & 深蓝 & \\
网格 & 深绿 & 每 64 像素一格，共 10 列 $\times$ 7 行\\
零轴 & 亮绿 & 屏幕垂直中心，对应 ADC 中点 512\\
CH0 波形 & 亮绿 & 默认通道\\
CH1 波形 & 黄色 & \\
文字 & 白色 & \\
\bottomrule
\end{tabular}
\end{center}

\subsection{底部文字读数}

文字区每 2 秒自动翻页，两页内容如下：

\begin{center}
\small
\begin{tabular}{lll}
\toprule
位置 & 第 1 页（时域页） & 第 2 页（测量页）\\
\midrule
第一行左 & \sig{AD=xxxx} 通道 0 当前值 & \sig{HI=xxxx} 最大值\\
第一行中 & \sig{U=xxxx} 通道 1 当前值 & \sig{LO=xxxx} 最小值\\
第一行右 & \sig{TB=xxxx} 采样分频 & \sig{AV=xxxx} 平均值\\
第二行左 & \sig{VP=xxxx} 峰峰值 & \sig{PE=xxxx} 周期\\
第二行右 & \sig{FR=xxxx} 频率 & \sig{DU=xxxx} 占空比\\
\bottomrule
\end{tabular}
\end{center}

\begin{tipbox}
进入按键菜单后，屏幕会\textbf{锁定在第 1 页}，并把第一行改为
\sig{AD=菜单项号}、\sig{U=当前值}，方便边调边看。退出菜单后恢复自动翻页。
\end{tipbox}

% ======================= 3 按键操作（核心） =======================
\section{按键操作}
\label{sec:key}

\subsection{按键位置}

板上共有 5 个按键，其中 \textbf{C 键为硬件复位}，其余 4 个为功能键：

\begin{center}
\small
\begin{tabular}{llll}
\toprule
按键 & 丝印 & FPGA 引脚 & 默认功能\\
\midrule
\textbf{U}（上） & BUTTON\_U & J13 & 垂直放大 / 进菜单\\
\textbf{D}（下） & BUTTON\_D & G12 & 垂直缩小 / 运行停止\\
\textbf{L}（左） & BUTTON\_L & H14 & 时基加快 / 通道循环\\
\textbf{R}（右） & BUTTON\_R & L12 & 时基减慢 / 波形切换\\
\textbf{C}（中） & BUTTON\_C & G11 & \textbf{复位}（不可重定义）\\
\bottomrule
\end{tabular}
\end{center}

\subsection{正常模式（默认）}

\textbf{短按}（按下后立即松开，$<0.8$\,s）：

\begin{center}
\small
\begin{tabular}{lll}
\toprule
按键 & 功能 & 效果\\
\midrule
U & 垂直放大一档 & 波形变高（档位 7$\rightarrow$2，最小 2 最高）\\
D & 垂直缩小一档 & 波形变矮\\
L & \textbf{水平放大}（1x$\rightarrow$2x$\rightarrow$4x） & \textbf{只放大已采数据，不重新采集}，画面不闪\\
R & \textbf{水平缩小}（4x$\rightarrow$2x$\rightarrow$1x$\rightarrow$全屏） & 同上，全屏档覆盖全部 1020 个采样点\\
\bottomrule
\end{tabular}
\end{center}

\textbf{长按}（按住超过 $0.8$\,s）：

\begin{center}
\small
\begin{tabular}{lll}
\toprule
按键 & 功能 & 效果\\
\midrule
U & 菜单 开 / 关 & 屏幕左上角出现或收起设置面板\\
D & 运行 / 冻结 & 冻结画面（STOP）或恢复采集（RUN）\\
L & \textbf{切换菜单页} & \textbf{仅在菜单已打开时有效}：示波器页 $\leftrightarrow$ 信号源页\\
R & 滚动模式 开 / 关 & 纸带式连续左移显示\\
\bottomrule
\end{tabular}
\end{center}

\begin{tipbox}
\textbf{按键无反应？} 请先确认是否处于菜单模式——菜单模式下短按 U/D 是"调整当前项"，
不再执行放大/缩小。长按 U 即可退出菜单。
\end{tipbox}

\subsection{菜单模式}

长按 U 进入菜单，屏幕第一行显示 \sig{AD=项号} 与 \sig{U=当前值}。

\begin{center}
\small
\begin{tabular}{ll}
\toprule
操作 & 作用\\
\midrule
短按 L / R & 左右移动光标（项号 0$\sim$7 循环）\\
短按 U / D & 当前项数值 加 / 减\\
长按 U & 退出菜单\\
\bottomrule
\end{tabular}
\end{center}

\subsection{菜单项一览}

菜单共两页，长按 L 切换（仅在菜单已打开时有效）。
\textbf{统一方向约定：U 键 = 减 / 放大 / 上移，D 键 = 增 / 缩小 / 下移。}

\paragraph{第 0 页：示波器设置（9 项）}
\begin{center}
\small
\begin{tabular}{clll}
\toprule
项号 & 功能 & 取值范围 & U / D 作用\\
\midrule
0 & CH0 开关 & 0 / 1 & 切换\\
1 & CH1 开关 & 0 / 1 & 切换\\
2 & CH0 垂直档位 (VS0) & 2$\sim$7 & U 放大 / D 缩小（数值越小波形越高）\\
3 & CH1 垂直档位 (VS1) & 2$\sim$7 & 同上\\
4 & CH0 垂直位置 (VP0) & 0$\sim$9，5 居中 & U 上移 / D 下移\\
5 & CH1 垂直位置 (VP1) & 0$\sim$9，5 居中 & 同上\\
6 & 水平缩放 (HSC) & 0$\sim$3 & 0=1x 1=2x 2=4x 3=全屏\\
7 & 时基 (TB) & 0$\sim$6 & U 更快 / D 更慢（\textbf{会重新采集}）\\
8 & 滚动模式 (ROL) & 0 / 1 & 开 / 关\\
\bottomrule
\end{tabular}
\end{center}

\paragraph{第 1 页：信号源设置（4 项，仅未接 ADC 时可用）}
\begin{center}
\small
\begin{tabular}{clll}
\toprule
项号 & 功能 & 取值范围 & U / D 作用\\
\midrule
0 & 波形 (WAV) & 0$\sim$3 & 0 正弦 1 方波 2 三角 3 锯齿\\
1 & 频率 (FRE) & 0$\sim$7 & 1 / 5 / 10 / 50 / 100 / 500 / 1k / 2k\,Hz\\
2 & 幅度 (ATT) & 0$\sim$3 & 100\% / 50\% / 25\% / 12.5\%\\
3 & 相位 (PHA) & 0$\sim$9 & 每档 36$^\circ$\\
\bottomrule
\end{tabular}
\end{center}

\begin{tipbox}
\textbf{接上 AD7606 后，第 1 页会整体消失}（面板只剩 9 行）。
此时信号来自外部 ADC，内置信号源的波形/频率/幅度/相位不再起作用。
\end{tipbox}

\begin{tipbox}
菜单项 6、7 只有在\textbf{未接 AD7606}（使用内置信号源）时才影响波形。
接入外部 ADC 后，波形来自 AD7606，这两项不再起作用。
\end{tipbox}

% ======================= 4 垂直与水平 =======================
\section{垂直系统与时基}

\subsection{垂直档位}

屏幕垂直方向共 448 像素。默认档位（4）对应"满量程铺满屏幕"：

\begin{center}
\small
\begin{tabular}{ccc}
\toprule
档位 & 相对增益 & 适用场景\\
\midrule
2 & $\times4$ & 小信号（只看波形顶部/底部细节）\\
3 & $\times2$ & 中等幅度信号\\
4 & $\times1$（默认） & 满量程显示\\
5 & $\times0.5$ & 大幅度信号\\
6 & $\times0.25$ & 接近满幅信号\\
7 & $\times0.125$ & 观察是否有削顶\\
\bottomrule
\end{tabular}
\end{center}

\subsection{垂直位置}

配合档位使用：先把波形缩小，再上下平移，即可把不同幅度的两路信号
\textbf{分开放置}，互不重叠。CH0 与 CH1 的位置独立可调。

\subsection{时基}

时基决定"屏幕上 640 个点覆盖多长时间"。7 个档位循环切换：

\begin{center}
\small
\begin{tabular}{cll}
\toprule
档位 & 采样分频 & 说明\\
\midrule
0 & 50000 & 最慢，适合看低频信号的全貌\\
1 & 10000 & \\
2 & 2000 & 演示用 50\,Hz 信号在这一档\\
3 & 500 & \\
4 & 100 & \\
5 & 10 & \\
6 & 1 & 最快，适合看边沿细节\\
\bottomrule
\end{tabular}
\end{center}

\begin{tipbox}
时基单位取决于 ADC 的实际采样率。接 AD7606 后采样率固定为 200\,kSa/s，
此时"档位 6"即每个像素一个真实采样点。
\end{tipbox}

% ======================= 5 函数发生器 =======================
\section{内置函数发生器}

未接 AD7606 时，内置信号源代替外部信号接入示波器，用于自检与演示。

\subsection{规格}

\begin{center}
\small
\begin{tabular}{ll}
\toprule
项目 & 参数\\
\midrule
采样率 & 12\,kSa/s（与 ADC 通路同接口）\\
波形 & 正弦、方波、三角波、锯齿波\\
频率 & 1 / 5 / 10 / 50 / 100 / 500 / 1000 / 2000\,Hz\\
幅度 & 100\% / 50\% / 25\% / 12.5\%（中心 512，满幅 $\pm400$）\\
分辨率 & 10\,bit（与 ADC 归一化后一致）\\
\bottomrule
\end{tabular}
\end{center}

\subsection{操作方法}

\begin{enumerate}[leftmargin=2em]
  \item \textbf{切换波形}：长按 \textbf{R} 键，循环切换四种波形
  \item \textbf{切换频率 / 幅度}：长按 \textbf{U} 进入菜单，用 L/R 选到第 6、7 项，
        再用 U/D 调整
  \item 波形立即生效，无需重启
\end{enumerate}

\begin{tipbox}
\textbf{为什么频率只有 8 档？}
档位是预先算好的相位增量表，切换时不会有任何跳变或毛刺。
后续版本会加入连续可调（用按键长按加速/减速）。
\end{tipbox}

% ======================= 6 触发系统 =======================
\section{触发系统}

\subsection{触发模式}

\begin{center}
\small
\begin{tabular}{lll}
\toprule
模式 & 行为 & 适用场景\\
\midrule
自动（AUTO） & 无触发也持续刷新 & 观察未知信号\\
正常（NORMAL） & 只有满足触发条件才刷新 & 观察稳定周期信号\\
单次（SINGLE） & 捕获一屏后冻结 & 抓取单次事件\\
\bottomrule
\end{tabular}
\end{center}

\subsection{触发参数}

\begin{center}
\small
\begin{tabular}{lll}
\toprule
参数 & 范围 & 说明\\
\midrule
触发源 & CH0 / CH1 & 用哪个通道做触发判断\\
触发边沿 & 上升 / 下降 & 有符号数在电平上下穿越时触发\\
触发电平 & 0$\sim$1023 & 默认 512（中点）\\
预触发点数 & 0$\sim$1023 & 触发点之前的采样点数\\
保持时间 & 0$\sim$65535 & 两次触发之间的最小间隔\\
\bottomrule
\end{tabular}
\end{center}

\begin{tipbox}
\textbf{画面一直不更新？} 检查是否处于 SINGLE 模式且已经捕获完成——
此时需要再次触发。短按 U/D 或按 D 长按（RUN/STOP）可重新采集。
\end{tipbox}

% ======================= 7 测量 =======================
\section{自动测量}

系统对每屏 1024 个采样点实时计算以下参数，显示在文字区（第 2 页）：

\begin{center}
\small
\begin{tabular}{lll}
\toprule
参数 & 符号 & 说明\\
\midrule
峰峰值 & \sig{VP} & 最大值 $-$ 最小值\\
最大值 & \sig{HI} & 一屏内最大采样值\\
最小值 & \sig{LO} & 一屏内最小采样值\\
平均值 & \sig{AV} & 一屏算术平均\\
周期 & \sig{PE} & 相邻两次过中点的上升沿间隔\\
频率 & \sig{FR} & 周期的倒数\\
占空比 & \sig{DU} & 高电平时间 / 周期\\
\bottomrule
\end{tabular}
\end{center}

% ======================= 8 串口 =======================
\section{串口控制}
\label{sec:uart}

除按键外，还可以用串口（115200 8N1）远程配置。命令格式为
\sig{<命令字母><十六进制参数><回车>}，例如 \sig{T1F40} 表示设置采样分频为 \sig{0x1F40}。

\begin{center}
\small
\begin{tabular}{cll}
\toprule
命令 & 功能 & 参数\\
\midrule
\sig{T} & 设置采样分频 & 16\,bit（\sig{0} = 不分频）\\
\sig{R} & 运行 / 停止 & \sig{0} / \sig{1}\\
\sig{A} & 触发使能 & \sig{0} / \sig{1}\\
\sig{L} & 触发电平 & 10\,bit\\
\sig{E} & 触发边沿 & \sig{0} 上升 / \sig{1} 下降\\
\sig{S} & 滚动模式 & \sig{0} / \sig{1}\\
\sig{M} & 触发模式 & \sig{0} 自动 / \sig{1} 正常 / \sig{2} 单次\\
\sig{C} & 触发源 & \sig{0} CH0 / \sig{1} CH1\\
\sig{P} & 预触发点数 & 10\,bit\\
\sig{H} & 保持时间 & 16\,bit\\
\bottomrule
\end{tabular}
\end{center}

\begin{tipbox}
串口下发的设置\textbf{优先于按键}。若发现按键改不动某个参数，
可能是串口刚下发过同一条命令。
\end{tipbox}

% ======================= 9 常见问题 =======================
\section{常见问题}

\begin{center}
\small
\begin{tabular}{p{5.2cm}p{9.2cm}}
\toprule
现象 & 处理办法\\
\midrule
屏幕全黑 & 检查 VGA 线是否插紧；确认 FPGA 已加载比特流（重新烧写或烧写到配置 Flash）\\
画面上下不一致 / 抖动 & \textbf{先怀疑显示器}。VGA 逐行扫描，同一列所有行读同一采样地址，
FPGA 不可能产生"上半正常、下半异常"的图像。重新插拔显示器电源；仍不行则换一台显示器\\
画面撕裂、同一条波形出现两份 & 无双缓冲导致的正常现象。长按 D 停止采集即可冻结画面\\
波形不动 & 处于 STOP 或 SINGLE 已捕获状态。长按 D 恢复运行\\
按键无反应 & 可能处于菜单模式（用长按 U 退出）；或按键被串口命令覆盖\\
波形成一条直线 & 未接 AD7606 时，检查菜单第 6/7 项是否把幅度设得太小；或时基档位不合适\\
频率读数乱跳 & 触发不稳定。调触发电平或改用自动模式\\
\bottomrule
\end{tabular}
\end{center}

% ======================= 附录 =======================
\appendix
\section{技术规格}

\begin{center}
\small
\begin{tabular}{ll}
\toprule
项目 & 参数\\
\midrule
FPGA & 安路 EG4S20BG256（19600 LUT，64 个 9K BRAM，4 个 PLL）\\
系统时钟 & 24\,MHz 晶振，PLL 倍频至 25.2\,MHz\\
VGA 时序 & 640$\times$480，像素时钟 25.2\,MHz，行频 31.5\,kHz，场频 60.00\,Hz\\
存储深度 & 1024 $\times$ 10\,bit 单口 RAM（每通道）\\
ADC & AD7606，16\,bit，200\,kSa/s，SPI 配置 + 并行读\\
按键 & 5 个（1 复位 + 4 功能）\\
显示资源占用 & LUT 约 2600 / 19600（13\%），BRAM 4 / 64（6\%）\\
时序余量 & 14.9\,ns（$@$25.2\,MHz）\\
\bottomrule
\end{tabular}
\end{center}

\section{版本历史}

\begin{center}
\small
\begin{tabular}{llp{8.4cm}}
\toprule
版本 & 日期 & 内容\\
\midrule
v1.0 & 2026-10-30 & 首版：4 键操作 + 菜单 + 内置函数发生器（4 波形 / 8 频率 / 4 幅度）\\
\bottomrule
\end{tabular}
\end{center}

\section{开发待完成清单}
\label{sec:todo}

本清单随固件迭代维护：$\checkmark$ 已完成，$\square$ 待完成。
优先级 P1 最高（影响可用性），P4 最低（锦上添花）。

\subsection*{已完成（截至固件 v1.1）}

\begin{center}
\small
\begin{tabular}{p{6.2cm}cp{6.0cm}}
\toprule
功能 & 状态 & 说明\\
\midrule
4 键操作（短按 / 长按） & $\checkmark$ & 垂直缩放、水平缩放、运行冻结、菜单\\
双页菜单 UI（P0 示波器 / P1 信号源） & $\checkmark$ & 黄框面板 + 选中行高亮\\
垂直档位与垂直位置 & $\checkmark$ & 6 档缩放，9 档位置\\
水平缩放（独立于时基） & $\checkmark$ & 1x / 2x / 4x / 全屏，不重新采集\\
时基 / 滚动 / 运行停止 & $\checkmark$ & 三者已按概念拆开\\
内置函数发生器 & $\checkmark$ & 4 波形 / 8 频率 / 4 幅度 / 10 档相位\\
ADC 模式下信号源自动禁用 & $\checkmark$ & \(\sig{GEN\_ENA=0}\) 时 P1 页整体隐藏\\
\bottomrule
\end{tabular}
\end{center}

\subsection*{待完成}

\begin{center}
\small
\begin{tabular}{clp{7.0cm}c}
\toprule
编号 & 功能 & 说明 & 优先级\\
\midrule
\multicolumn{4}{l}{\textbf{A. 交互与显示}}\\
\midrule
A1 & 触发电平虚线 & 屏幕上用虚线标出当前触发电平位置，调电平时可直观看到 & P2\\
A2 & 触发点标记 & 在触发位置画竖线或三角标记，配合预触发可看出触发点 & P2\\
A3 & 采集状态指示 & 屏幕左上角显示 ARMED / TRIG'D / STOP / AUTO & P2\\
A4 & 光标测量 & 两条可移动游标，显示 $\Delta T$ 与 $\Delta V$ & P4\\
A5 & 通道颜色可选 & 绿 / 黄 / 青 / 品红，便于双通道区分 & P4\\
A6 & 菜单交互增强 & 选中项闪烁、长按加速调节 & P4\\
\midrule
\multicolumn{4}{l}{\textbf{B. 采集与信号}}\\
\midrule
B1 & AUTOSET 自动量程 & 用已有 Vpp / 频率反推合适的垂直档位与时基，一键出图 & P3\\
B2 & 双通道真并行采集 & 当前 4 通道交错写入同一 RAM，属设计缺陷，需每通道独立存储 & P3\\
B3 & 双缓冲显示 & 消除跟踪模式下的撕裂，需 SDRAM 或乒乓半缓冲 & P4\\
B4 & Vrms 有效值测量 & 增加平方累加电路 & P3\\
B5 & 频率连续可调 & 目前仅 8 档，需连续相位增量 & P3\\
B6 & 外部 ADC 实测 & 接 AD7606 后跑通全流程并验证量程/线性度 & P2\\
B7 & 预触发 / 保持时间界面化 & 目前只能通过串口设置 & P3\\
\midrule
\multicolumn{4}{l}{\textbf{C. 板载资源利用}}\\
\midrule
C1 & LED 状态指示 & LED[3:0] = RUN / ROLL / TRIG / DONE；LED[4] 菜单；LED[5] PLL 锁定；LED[15:8] 作为当前采样值的电平条 & P2\\
C2 & 数码管显示信号源参数 & 显示波形编号 / 频率档 / 相位档，便于盲操作 & P2\\
C3 & 拨码开关扩展 & 若板上有拨码开关，可用于模式选择（如正常/调试） & P4\\
\midrule
\multicolumn{4}{l}{\textbf{D. 串口与通信}}\\
\midrule
D1 & 新增参数命令 & 现有 10 条命令尚未覆盖水平缩放、滚动、信号源参数 & P3\\
D2 & 设置导出 / 状态回读 & 上位机读取当前全部配置 & P4\\
\midrule
\multicolumn{4}{l}{\textbf{E. 交付与文档}}\\
\midrule
E1 & 配置 Flash 固化 & 让设计断电不丢失，上电自动加载 & P1\\
E2 & 手册随固件同步 & 每次固件更新后同步本手册与回归说明 & P1\\
E3 & 按键行为回归用例 & 目前仿真是数据通路级，缺按键/菜单行为测试 & P2\\
\bottomrule
\end{tabular}
\end{center}

\begin{tipbox}
\textbf{B2（双通道并行采集）与 B3（双缓冲）是架构级改动}，
需要修改存储结构并重新综合时序，建议在功能稳定后单独安排。
其余项均为增量修改，可按优先级逐个交付。
\end{tipbox}

\end{document}